\noindent You are a Research Scientist in the early brainstorming phase of a project. Your task is to write a high-level "Concept Note" (\texttt{idea.md}) based on a provided text.

\vspace{0.5em}
\noindent You have been given the text content of a paper ([PAPER CONTENT]). You must distill this into a streamlined, conceptual project proposal.

\vspace{1em}
\noindent\textbf{Critical Data Ingestion Rules}
\begin{enumerate}
    \item \textbf{Target Content:} Extract information \textit{only} regarding the \textbf{concept and intuition} of the research.
    \begin{itemize}
        \item \textbf{Focus Areas:} Problem Definition, Motivation, High-level Method/Algorithm.
    \end{itemize}
    \item \textbf{Exclusion Zone:} Stop extracting information once the text shifts to \textbf{empirical verification}.
    \begin{itemize}
        \item \textbf{STRICTLY IGNORE:} Experiments, Results, Evaluation, Comparisons, Ablations, or Conclusions.
    \end{itemize}
\end{enumerate}

\vspace{1em}
\noindent\textbf{Instructions}
\begin{enumerate}
    \item \textbf{Perspective:} Write in \textbf{First-Person Future Tense} (e.g., "We propose to explore...", "We aim to investigate...").
    \item \textbf{Enforce Sparsity (High-Level Focus):}
    \begin{itemize}
        \item \textbf{Conceptual Over Mathematical:} \textbf{Do NOT use LaTeX.} Do not provide formulas. Instead of writing the math, describe the \textit{intuition} or \textit{purpose} of the component (e.g., "We will use a loss function designed to maximize perceptual similarity...").
        \item \textbf{Strategic Logic:} Describe the methodology at a "whiteboard" level. Avoid hyperparameters (like "\$d=512\$"). Focus on the flow of data and the logic of the modules.
        \item \textbf{Simulation:} Mimic the early design phase where the intuition is clear, but the exact implementation details are not yet finalized.
    \end{itemize}
    \item \textbf{Structure:}
    \begin{itemize}
        \item \textbf{Problem Statement:} The gap we are filling.
        \item \textbf{Core Hypothesis:} The specific technical novelty.
        \item \textbf{Proposed Methodology:} A conceptual description. Focus on strategy and logical steps. Describe modules by their function, not their math.
        \item \textbf{Expected Contribution:} The theoretical value.
    \end{itemize}
    \item \textbf{Formatting:}
    \begin{itemize}
        \item Be self-contained. \textbf{No citations, no URLs, no references to Figure/Table numbers.}
        \item Fully anonymize authors/titles.
    \end{itemize}
\end{enumerate}

\vspace{1em}
\noindent\textbf{Output Format}\\
\vspace{0.5em}
Return \textit{only} the markdown memo in the following structure:

\begin{verbatim}
```markdown
## Problem Statement
(Precise definition of the technical problem.)

## Core Hypothesis
(The proposed solution/intuition.)

## Proposed Methodology (High-Level Technical Approach)
(A conceptual description of the approach. Focus on the strategy and 
logical steps rather than mathematical derivations. Describe the modules 
and their functions.)

## Expected Contribution
(The intended theoretical or practical value of this approach.)
```

[PAPER CONTENT]
{paper_content}
[END PAPER CONTENT]
\end{verbatim}